\documentclass[11pt,letterpaper,oneside,reqno]{amsart}
\usepackage[letterpaper,textwidth=6.15in,textheight=8.82in,footskip=28pt,centering]{geometry}
\usepackage{amssymb,mathtools,booktabs,longtable,array,enumitem,needspace}
\usepackage[final]{microtype}
\usepackage[hidelinks]{hyperref}
\pagestyle{plain}
\setlength{\emergencystretch}{1.5em}
\setlist{leftmargin=*,nosep}
\raggedbottom
\newtheorem{theorem}{Theorem}
\newtheorem{lemma}[theorem]{Lemma}
\newtheorem{prop}[theorem]{Proposition}
\newcommand{\T}{\mathcal T_S}
\newcommand{\N}{\mathbb N}
\newcommand{\eps}{\varepsilon}
\newcommand{\pair}[2]{\langle #1,#2\rangle}
\newcommand{\word}{\mathsf W}
\newcommand{\bx}[1]{\mathsf{box}(#1)}
\newcommand{\slot}{\mathsf{slot}}
\newcommand{\live}{\mathsf{live}}
\newcommand{\cmt}{\mathsf c}
\newcommand{\pops}{\mathsf d}
\newcommand{\st}{\mathsf{stable}}
\newcommand{\arity}{\ell}
\newcommand{\fired}{\mathsf{fired}}
\newcommand{\halt}{\mathsf{halt}}
\newcommand{\bin}{\{0,1\}}
\newcommand{\Rule}{\longrightarrow_S}
\newcommand{\GH}{https://github.com/Bobtheotherone/pure-s-combinator-research}
\DeclareRobustCommand{\PaperAuthors}{Radames Mercado-Barbosa and Anthony Langdon}

\hypersetup{pdftitle={Computational universality of the pure S-calculus under a uniform selected reduction},pdfauthor={\PaperAuthors}}
\title[Selected reduction in the pure \(S\)-calculus]
  {Computational universality of the pure \(S\)-calculus under a uniform selected reduction}
\author{\PaperAuthors}
\date{26 November 2026}
\begin{document}
\begin{abstract}
We give a single reduction strategy for the pure \(S\)-calculus that simulates every finite binary cyclic tag program. Each nonempty tag step takes two, three, or four contractions of \(Sxyz\to(xz)(yz)\). A contraction consumes a queue symbol, and the word released by a later contraction becomes part of the same queue. For a fixed universal program, a regular tree language detects halting and the terminal term determines the numerical output. We prove the input/output equivalence in Lean~4, both at tag-step boundaries and for an eight-mode controller that contracts one redex at a time. To bound the power of the auxiliary operations, we also prove a finite cutoff for acceptance in systems that combine whole tree registers under finite observation.
\end{abstract}
\maketitle

\section{Introduction}
The \(S\) rule duplicates its third argument and retains all three arguments. To represent a queue in such a system, one must distinguish the part still to be read from the copies left by earlier contractions. Our encoding keeps the unread letters in order. After a letter is consumed, its residual records the bit that determines which program word to release.

Let \(\T\) be the finite ordered binary trees with leaf \(S\) and application as internal node. Juxtaposition associates to the left. A contraction replaces one occurrence, even when other occurrences are syntactically equal. A \emph{regular redex selector} is a fixed regular language of singly marked finite trees that marks exactly one legal redex on every reducible input. It gives a function on bare terms. On normal terms we extend the function by the identity.

A cyclic tag program \(P=(v_0,\ldots,v_{p-1})\), \(p\ge1\), consists of binary words, possibly empty. In phase \(i\), a source step is
\begin{equation}
  (i,bz)\longmapsto(i+1\bmod p,\ z v_i^{\,b}),\qquad
  v^0=\eps,\quad v^1=v.                                      \label{cts}
\end{equation}
The empty queue ends the run. Cook proves universality of this model by simulation of tag systems \cite{cook}. We emulate its whole queue, for every program and word.

\begin{theorem}\label{main}
There are a linear-size constructor encoder \(E(P,w)\) and a single root-restarted, finite-state regular selector \(\sigma\), independent of \(P\), with the following property. Starting with phase zero and queue \(w\), every nonempty step \eqref{cts} is realized between consecutive distinguished samples of the native path, at distance \(2\) to \(4\). The live words at these samples are exactly the source queues. A fixed regular test detects an empty queue precisely at completion, including initial emptiness. Thereafter the native path can continue indefinitely with the same decoded terminal queue.
\end{theorem}

Wolfram allows selected reduction and output detection along a nonterminating path \cite{announcement,prize}. The bound counts contractions. Search and copying costs depend on the term representation. Theorem~\ref{output} reads a register from the live word. Theorem~\ref{certnumeric}, formalized in Lean, reads a unary numeral from the halted tape.

\section{Encoding the queue and the program}
For terms \(u,v\), write \(\pair u v=Suv\). Let \(\mathsf{nil}\) denote the empty list and \(u::\ell\) the list with head \(u\) and tail \(\ell\). The encoding of lists of terms is given by
\[
  \operatorname{enc}(\mathsf{nil})=S,\qquad
  \operatorname{enc}(u::\ell)=\pair{u}{\operatorname{enc}(\ell)}.
\]
We enclose each visible word \(X\) in the term \(\bx{X}=\pair S X\). In an address, \(L\) and \(R\) denote a step to the left and right child, respectively. The fields of a pair have addresses \(LR\) and \(R\). A contraction acts at the selected occurrence, even when the same subtree occurs elsewhere in the term.

For the two bit values, set \(B_0=S,\ B_1=SS,\ J_b=SSB_b\), and put
\[
 \word(\eps)=S,\qquad \word(bw)=J_b\word(w).
\]
The read operation is a single contraction
\begin{equation}
 J_bX\ \Rule\ (SX)(B_bX).                                    \label{read}
\end{equation}
Call the source a live cell and the result a spent cell. In the latter, \(LR\) continues the word and \(RL\) retains the consumed bit. The rightmost copy is never touched. After subsequent contractions the spent shape is \((SA)(B_bG)\) with arbitrary history \(G\). For a live cell the continuation is at \(R\). A stream of such cells records all bits in their original order and records which of them have been consumed. Its terminal is \(S\). Equation~\eqref{read} and this history convention occur in the prior work described in Section~\ref{sources}.

To store a program word until it is called for, define \(K(V)=SSSV\). The two arms of a slot are
\begin{equation}
 \slot(v)=\pair{K(S)}{K(\bx{\word(v)})},\qquad
 K(V)\ \Rule\ (SV)(SV).                                      \label{commit}
\end{equation}
On input bit \(b\), fire arm \(b\), with the first arm numbered zero. A fired arm exposes \(V\) at \(LR\). For arm one, its right child is the appended word. For arm zero the appended word is empty. A fired and subsequently read arm has shape \((SV')G\). Its right field and the unused arm remain opaque. No equality of copies is used to recognize a fired arm. In particular, a fired word can be consumed without changing the master copy of the program.

For a program of period \(p\), the initial term contains one list of slots. Its other components are
\begin{align}
 L&=\operatorname{enc}\bigl((\slot(v_i))_{i<p}\bigr),&
 F&=\pair{S}{L},& A&=SF, \notag\\
 X_0&=SAA,& X_{j+1}&=A X_j,&
 E(P,w)&=\pair{\bx{\word(w)}}{X_0X_0}.                             \label{seed}
\end{align}
Only \(X_0\) is needed by the encoder. The descendants called \(X_j\) arise during reduction. A region for program periods has a terminal pair, called its frontier. It starts at \(X_0X_0\). Two root contractions describe its complete allocation mechanism:
\begin{align}
 X_0X_j &\Rule (A X_j)(A X_j)=X_{j+1}X_{j+1},                 \label{reset}\\
 X_{r+1}X_j &\Rule (F X_j)(X_rX_j).                           \label{alloc}
\end{align}
In \eqref{alloc}, the left child contains a pristine period. Its list is at \(LLR\) from the new root. The right child is the new frontier. Equation~\eqref{reset} is followed by an instance of \eqref{alloc} with \(r=j\) when a period is needed. A fresh period therefore costs one or two contractions. All previously allocated periods lie outside this occurrence. The specialization \(A=SF\) adapts the growable-register allocator of~\cite{five} to whole strings of executable slots.

\section{The selected step}
At any time the second field of the seed is a right spine of allocated frames
\((F'X_j)t\), followed by a frontier. \(F'\) is a copy of \(F\) with some slots fired and some released words subsequently consumed. Neither the \(X_j\) field of a frame nor any subtree of a frontier is separately edited. Read frames from the outside inward, and slots inside each frame in list order.

Form the ordered sequence of \emph{visible words} from the header followed by the released words of all fired one-arms in slot order. A fired zero-arm contributes the empty word. Walk just the designated continuations of these words. Write \(\cmt\) for the number of fired slots, \(\pops\) for the number of spent cells in these visible words, and \(V\) for the concatenation of their live letters. These counts are proof quantities. The policy tests their parities. It uses the last spent bit only when such a bit exists.

\begin{center}
\small
\begin{tabular}{@{}p{.43\linewidth}p{.50\linewidth}@{}}
\toprule
Condition on the current bare tree & Occurrence selected for one contraction\\
\midrule
\(\pops\equiv\cmt\pmod2\) and \(V\ne\eps\)
 & The first live cell of the first visible word that still has a live cell.\\[3pt]
\(\pops\not\equiv\cmt\pmod2\) and an unused slot exists
 & Arm \(b\) of the first unused slot, where \(b\) is the last spent bit.\\[3pt]
\(\pops\not\equiv\cmt\pmod2\) and no unused slot exists
 & The frontier, by \eqref{reset} or \eqref{alloc}.\\[3pt]
\(\pops\equiv\cmt\pmod2\) and \(V=\eps\)
 & The frontier. Completion persists while unused periods accumulate.\\
\bottomrule
\end{tabular}
\end{center}
The phase is also structural: in the block containing the first unused slot, count its preceding fired slots. If no unused slot exists, the next phase is zero. This gives \(\cmt\bmod p\) even when empty post-halt blocks accumulate.

A sample is \(\st\) when the parities agree on the parsed representation. Section~\ref{regular} defines the policy away from this representation.

\begin{proof}[Proof of Theorem~\ref{main}]
Initially no slot is fired and the sole visible word is \(w\). We maintain three facts. Fired slots form an initial segment of the allocated slots. Spent letters form an initial segment of all visible letters. Finally, \(\pops=\cmt\) at a stable sample and \(\pops=\cmt+1\) between active stable samples.

Suppose a stable sample has \(V=bz\) and \(\cmt=n\). Equation~\eqref{read} spends exactly its first live letter. Both residuals of the tail are present, but only one remains visible. The concatenated live word is now \(z\), and the last spent tag is \(b\). If needed, \eqref{reset}--\eqref{alloc} supply the next slot. They leave both counts and all visible words unchanged. Because the slots spell successive identical periods and the first \(n\) are fired, the next one is \(\slot(v_{n\bmod p})\). Equation~\eqref{commit} fires its selected arm and adds exactly \(v_{n\bmod p}^{\,b}\) after every earlier visible letter. Thus the new \(V\) is \(zv_{n\bmod p}^{\,b}\) and exactly one additional slot is fired. All three facts hold again. None of the intervening samples has equal parities. There was one pop, one firing and zero to two allocations.

This proves the assertion for arbitrarily many consecutive steps by ordinary induction. At a stable empty queue the last policy row preserves emptiness and both counts forever. The frontier always has a native redex. Every encoded path consequently contains only legal contractions even after the source finishes.
\end{proof}

\textit{Example.} With \(P=(1)\), \(w=1\), step one exhausts the header and releases a single \(1\). Step two must consume the released occurrence. Inductively every subsequent bit comes from the immediately preceding firing. Changing the sole appendant to \(0\) produces \(1,0,\eps\) at successive stable samples with exactly the same policy.

\section{Regularity and total definitions}\label{regular}
For a ground tree let \(\arity(S)=0\) and \(\arity(uv)=1+\arity(u)\). Tests of \(\arity=k\) for fixed \(k\) are bounded-depth tests. A frame has \(\arity=4\), a frontier \(\arity=3\). A slot arm before firing has head arity three and after firing has head arity two. A live cell's left child has head arity two and a spent cell's left child has head arity one. These distinctions survive every selected edit to a released payload.

Here are precise parsing conventions for arbitrary trees. Require the outside pair, a finite right spine of frames with left field \(F'X\), \(F'=\pair{S}{L'}\), and a frontier of arity three. Each \(L'\) is a nonempty list of slots. The field \(X\) and the children of the frontier are opaque. Each unused arm is \(K(G)\); each fired arm has shape \((SA)G\), with \(G\) arbitrary. A slot has at most one fired arm. Require the payload \(A\) of a fired one-arm and the input header to be boxes, and inspect their right children as words under the following grammar, with \(G\in\T\):
\[
 W ::= S\ \mid\ J_b W\ \mid\ (SW)(B_bG).
\]
The zero-arm payload may be ignored. Require fired slots to be a prefix and spent letters to be a prefix, in the orders of Section 3. These conditions over-approximate reachable terms. No test of pristine-copy equality, index equality, slot periodicity, or full history correctness is necessary.

\begin{lemma}\label{regularlem}
The parse, the stable and empty tests, and the graph of the designated marked occurrence are regular tree languages. The completed selector is a total function on \(\T\).
\end{lemma}
\begin{proof}
All productions above have fixed finite constructor patterns and disjoint alternatives at each role. The parser can send ignored fields to the unrestricted nonterminal \(G\). Extract a word of finite tokens in order: header-cell tokens first, then for each slot a fired/unused token and its visible-cell tokens. Each cell token reports live/spent and its finite letter. This traversal is a finite collection of recursive constructor clauses. Each clause emits fixed tokens and concatenates the streams of the designated children, each used at most once. For any word DFA, a bottom-up tree automaton stores, for each parsing role, the partial transformation on DFA states induced by the resulting stream. A failed clause contributes failure. Composition in the specified order realizes concatenation. Lists, frames, cells and the top pair require only finitely many roles. This constructs the automaton without materializing the stream. The two prefix conditions, parities, emptiness and the last spent bit are DFA tests.

On a singly marked tree, attach to a token the flag that the designated root lies at its pop, arm, or frontier address. Emit a final frontier token. Local fixed address tests and a three-valued mark count (zero, one, multiple) suffice. A word DFA tests that the mark designates the first eligible live cell or unused slot, the correct arm, or the frontier as prescribed. A product with the parity and tag tests settles all policy cases. The relation is functional on its projection, and each marked position is additionally required to be a native redex. Where the parse or nomination fails, choose the first redex in ordinary preorder. Its marked relation and the absence of any successful nomination are regular, the latter by projection and complement. On normal terms return the term itself. The actual encoded paths always lie in the proved cases.
\end{proof}

Thus one may implement selection by a fixed bottom-up pass and a descent to the unique marked position, recomputing annotations on the current tree. For example, the bottom-up pass stores the unmarked state and the finite set of possible states with one marked descendant. A downward pass carries the accepting context states and follows the unique possible witness. With an explicit occurrence tree this takes \(O(|t|)\) time and \(O(|t|)\) temporary finite annotations. The following construction removes the need to store these annotations.

\begin{prop}[A root-restarted tree walker]\label{walker}
On binary cyclic programs the selector can start every invocation at the root in the same finite control, observe only node kind and incoming side, and move one edge at a time. It takes \(O(|t|+1)\) elementary moves on every bare term. It has no persistent cursor, marks or work tape.
\end{prop}
\begin{proof}
One read-only scan computes the two parity bits, the last spent bit and flags for any live cell and unused slot. It checks the prefix conditions as it goes. A second scan from the root stops at the prescribed occurrence. It carries only the mode and, for a commit, the one payload-selection bit. All fixed-depth tests in either scan are implemented by a bounded sequence of child/parent moves, returning to the testing position.

The following return paths specify the part of this traversal where an implicit stack would otherwise be a concern.
\begin{enumerate}
\item From a word box descend \(R\). Walk down \(R\) at a live cell and \(LR\) at a spent cell. To return from the terminal \(S\), go up one edge at a time from a word node. A parent of arity three is the previous live cell. A parent of arity one has the form \(SA\); ascend once more to the previous spent cell. A parent of arity two is the box. These cases are disjoint, including for an empty word. The paths are inverses of the checked descent.
\item The header box is at \(LR\) of the whole term. In a fired one-arm the box is at \(LR\) of that arm, which is at \(R\) of the slot. Each return to the slot or to the header therefore takes a fixed number of moves. At a list cell enter its slot by \(LR\), return by two upward moves, and continue the list by \(R\).
\item At the list terminator ascend through list cells. Each is a pair whose first field is a slot, hence an application. The first parent pair with first field \(S\) is the template \(F'\). Ascend twice to the frame and take \(R\) to the next frame or frontier. Tests for a frame, a frontier and a genuine slot use bounded depth. Dormant program words and copied history are skipped.
\end{enumerate}
On a failed check, climb to the root and do a preorder search for the first actual redex. Incoming side implements entry and the two return modes for this depth-first search. If none exists the term is normal. On malformed input, every completed return above retraces a previously verified downward segment, and a failed probe goes straight to fallback. Each of the two scans visits disjoint word/list/frame regions in forward order, with only a bounded number of visits per node and bounded probes per visit. Thus the running time is linear on all inputs. The second scan changes the tree only at its final verified contraction. This gives the same choice on the promised terms and the regular fallback of Lemma~\ref{regularlem} when a required test fails.
\end{proof}
For a fixed finite queue table in the next section, decoding a finite tag, checking the menu layout and returning from an arm take a bounded number of additional movements depending only on its finite dimensions. The word boxes and period-list traversal are unchanged. The same two-pass implementation therefore applies to that fixed specialization.

\section{Finite menus and register output}\label{numeric}
The allocation and word operations also simulate a finite queue machine. The queue alphabet can then contain a letter for each register, so that its multiplicity gives the output.

A deterministic finite queue machine has finite states \(Q\), finite alphabet \(\Sigma\), and nonterminal transitions
\[
       (q,a z)\longmapsto(q',zv),\qquad T[q,a]=(q',v).
\]
Choose fixed orders on \(Q,\Sigma\), and encode their indices by \(B_0=S,\ B_{i+1}=S B_i\). Equation~\eqref{read} applies verbatim to any fixed finite alphabet. Use one slot per allocated period. This slot is a binary menu of arms indexed in row-major order by \((q,a)\). Its arm is
\begin{equation}
                     K\bigl(\pair{B_{q'}}{\bx{\word(v)}}\bigr).       \label{menuarm}
\end{equation}
Here the subscript \(q'\) denotes the fixed index of the state. For an explicit layout, pad the row-major list to the next power of two by repeating the last arm. A leaf containing \(a'\) is \(\pair{SS}{a'}\); an internal node with children \(l,r\) is \(\pair{S}{\pair l r}\). No padding leaf is selected. Use arbitrary dummy rows for a terminal state. The fixed template and seed are \eqref{seed} with \(L=\operatorname{enc}(\mathrm{menu}::\mathsf{nil})\). The same frontier equations apply.

The current state is the tag of the last fired arm, or \(q_0\) if none has fired. Each fired arm exposes its state and its word at \(LR\); the word within that pair is at \(RR\). Read all such words in order, after the header. The policy of Section 3 now fires the arm indexed by the current state and the last spent letter. At a stable terminal state it takes only the frontier action. The state transition and the word in \eqref{menuarm} are read from the released payload. The selector contains the finite key indexing convention and the halt-state test. Changing every target state and every appendant in a table of the same dimensions leaves the selector valid without modification.

\begin{lemma}\label{queue}
At stable samples the decoded state and queue of this construction realize every successive finite queue-machine transition on a nonempty queue. Each such transition costs two to four contractions. A stable terminal state persists with its output.
\end{lemma}
\begin{proof}
After popping \(a\), the previous state is still \(q\). Allocation changes neither datum. Firing \((q,a)\) releases one copy of \(\pair{B_{q'}}{\bx{\word(v)}}\) after all older active words. Its tag becomes the new current state and its letters give exactly \(zv\). The prefix and count invariants from Theorem~\ref{main} give induction and persist during terminal allocation. For regularity, a finite menu parser records whether zero or one arm is fired and, if appropriate, its state. The finite-state summaries of Lemma~\ref{regularlem} now track the last released state and the marked arm's key. A malformed menu falls under the same fallback.
\end{proof}

\paragraph{\textit{Compiling registers into the fixed table.}}
A register program with \(r\ge1\) registers operates by increment with a target, conditional decrement with two targets, unconditional jump, and halt. Use two different queue letters \(a_i,\#_i\) per register. At each instruction boundary its queue is the literal row
\begin{equation}
       R(n_0,\ldots,n_{r-1})=
       a_0^{n_0}\#_0\,a_1^{n_1}\#_1\cdots a_{r-1}^{n_{r-1}}\#_{r-1}. \label{row}
\end{equation}
States are \((q,f)\) with one flag \(f\in\{0,1\}\), initially zero. During one sweep pop each old letter and append a copy according to the following rules. Increment \(i\) replaces the output for \(\#_i\) by \(a_i\#_i\). Decrement \(i\) suppresses the first \(a_i\), setting \(f=1\), and copies any subsequent \(a_i\). Jump copies every letter. On reading the final delimiter choose the prescribed successor, including the zero or positive successor according to the accumulated flag, and clear \(f\). When incrementing the last register, emit its extra letter before changing state. All of these instructions specify individual entries of the finite table \eqref{menuarm}.

\section{The register-machine encoding}
Every old letter precedes every newly appended letter until the old final delimiter is read. Consequently one complete sweep produces exactly the row \eqref{row} for the instructed successor. The flag in a decrement sweep is one precisely when an old unit of the selected register was present. The last delimiter occurs once per row, including when every register is zero. Induction on instructions now proves the simulation for every finite run. The queue is nonempty at each sweep boundary. Before the final delimiter of a sweep is consumed, an old delimiter remains. The intermediate empty word between pop and commit cannot trigger a stable halt.

Fix the interpreter \(U\) in appendix A of the observed-computation manuscript \cite{observed}. It has 446 primitive labels and 29 registers. Appendix~\ref{app:code} lists its instructions. Its registers \(E,X,\mathit{program},\mathit{argument}\) are initialized to \(e,x,e,x\), in that order, and all others to zero. Its output register has zero-based index 4. Write \(\nu(e,x)\) for this vector, and \(T_U\) for its queue table constructed in Section~\ref{numeric}. Its finite code is fixed independently of \(e,x\). We may now take
\begin{align}
 E_U(e,x)&=\pair{\bx{\word(R(\nu(e,x)))}}{X_0 X_0}\quad
        \text{with the fixed template for }T_U, \notag\\
 \sigma_U&=\text{the completed one-contraction menu selector of Lemma \ref{queue}},\notag\\
 H_U(t)&=\mathbf 1\{\text{valid parse, stable parities, last state is terminal}\}, \notag\\
 D_U(t)&=
 \begin{cases}
    |V(t)|_{a_4}, & H_U(t)=1,\\
    \bot, & H_U(t)=0 .
 \end{cases}                                                     \label{maps}
\end{align}
The initial state convention applies when there are no fired slots. For the specified \(U\), the unique reachable terminal state is the halt label with cleared flag. One can designate just that state for the detector and policy.

\begin{theorem}\label{output}
For all \(e,x,y\in\N\),
\[
 U(e,x)\downarrow=y\ \iff\
 \exists k\ge0\ \bigl(H_U(\sigma_U^k E_U(e,x))=1\ \land\
                      D_U(\sigma_U^k E_U(e,x))=y\bigr).
\]
\end{theorem}
\begin{proof}
A finite interpreter run has finitely many finite sweeps. Lemma~\ref{queue} realizes each of their letters and enters the terminal state on the final delimiter of the last sweep. The row at that time contains exactly \(y\) live \(a_4\)'s. In the reverse direction a stable terminal sample can first occur only at initialization or on committing an actual transition into the terminal state. That commit closes a complete sweep. The sweep induction identifies a genuine halted interpreter configuration. All subsequent contractions preserve its row. This also proves uniqueness of every accepted numerical answer.
\end{proof}

\paragraph{\textit{Universality of the register model.}}
The original indexing uses \(0\) for the empty list and \([u\mid v]=2^u(2v+1)\) for a cell. Instruction zero halts; increment is \([0\mid[j\mid l]]\); conditional decrement is \([1\mid[j\mid[l_0\mid l_1]]]\). The interpreter fetches the instruction, splits its finite fields by repeated halving, traverses the encoded register list, updates its selected entry and reconstructs the list. Malformed code enters a self-loop. Its finite expansion uses unit increments, guarded decrements and jumps. Each valid list operation terminates and preserves the represented values. Independently, two natural registers can encode the strings to the left and right of a tape head as base-\(B\) stacks with nonzero symbol digits. Finite register programs implement push by fixed-base multiplication and addition, and pop by fixed-base quotient and remainder, using scratch registers. The scanned symbol and tape-machine state require finite control. Every tape instruction therefore has a terminating register subroutine, including extension by a blank. This gives a program for every partial recursive numerical function \cite{register}. Our \(S\)-encoder prints the two unary inputs and fixed syntax; it never runs the interpreter.

\section{A certified evaluation on bare terms}
\label{certified}
The same seed, slot arms and allocator admit a second evaluation policy. It acts on bare ground terms, with observations taken between successive tag steps. Between observations it carries one bit and a finite control mode. We describe the boundary map first and then its realization by individual contractions.

Use the following traversals with immutable input. On a word, a pop contracts its first live cell, descending through spent cells at \(LR\). On a slot it searches only the word in a fired one-arm. On a list of slots it tries them in order. On the rail it tries each allocated period in order and skips the frontier. On the complete tree it first tries the header and then the rail. A firing traverses just the rail in the same order. At the first slot with two unfired arms it contracts the arm designated by its bit parameter. An allocation skips the frames and contracts the frontier. At every success the operation rebuilds only the ancestors of the contracted occurrence. On a failed subsearch it continues with the next listed subsearch. The functions are total structural recursions on finite trees. The boundary operation is given by
\[
\begin{array}{ll}
 b\gets h(t).& b=\bot\ \Longrightarrow\ \mathrm{fail}.\\
 t_1\gets\mathrm{pop}(t).&\\
 \text{For }j=0,1,2:&
  \mathrm{fire}_b(t_{j+1})=u\ \Longrightarrow\ \mathrm{return}\ (2+j,u).\\
 &j<2\ \Longrightarrow\ t_{j+2}\gets\mathrm{allocate}(t_{j+1}).
\end{array}
\]
The partial action returns failure when a necessary operation fails. Its total boundary map \(B\) keeps the initial term on failure. A separate verified single-contraction realization has eight modes: boundary, three searches each with one bit, and a terminal fallback. Every invocation begins navigation at the root. The first search fires the chosen arm if present, or allocates and passes to the second search; the second does likewise; the third can fire. Only the bounded stage and the borrowed bit can pass to the next invocation. On an unavailable action the controller tries ordinary preorder, enters the fallback mode on success and reports normality precisely when there is no redex. Lean proves this normality equivalence and one-step legality for every finite input and every mode. For a successful boundary action of reported length \(k\), precisely \(k\) calls take the controller from boundary to boundary at exactly the specified whole ground tree. On the encoded domain every call on a nonempty queue succeeds. An empty queue has already been detected. There is no call to a source interpreter. A source-step counter is used only to state the correctness theorem.

Here \(h(t)\in\{\bot,0,1\}\) is a finite tree-algebra value. On every reachable boundary term it is the first live symbol or \(\bot\). A word observer on \(S\) gives \(\bot\), gives \(b\) on \(J_b w\), and on a spent cell follows the designated left continuation. Combine ordered fields by \(a\triangleright b=a\) if \(a\ne\bot\), and \(b\) otherwise. On a slot look only at the first word inside its fired right arm. Combine these observers over lists and frames, with a final first-choice between header and rail. Constant-depth shape tests give a total interpretation on arbitrary trees. The actual bottom-up state has five Boolean and fifteen ternary fields, including the promised head of a word, a slot, a row, a rail and the whole tree. The other fields store bounded-depth pattern matches on the two children. The finite transition table, expressed as constructor equations, is also given in the symbolic supplement. Lean checks an injection of this state into twenty finite coordinates. The detector is the single fixed regular language \(H(t)=[h(t)=\bot]\).

\begin{theorem}\label{certmain}
Let \(B\) be the bare boundary function and \(V\) the structural live-word projection. For every nonempty finite binary cyclic program, every initial word and every \(n\), \(V(B^nE(P,w))\) is the source queue after \(n\) absorbing source steps. On every finite tree, every successful action is a chain of exactly its reported \(k\) contextual \(S\) contractions, where \(2\le k\le4\). For the explicit fixed tape-to-Cook compiler \(C\), fixed program \(P_*\), and every finite deterministic binary tape-machine instance \(M\),
\[
             M\downarrow\quad\Longleftrightarrow\quad
              \exists n\in\mathbb N:\ H\!\left(B^n E(P_*,C(M))\right)=1.
\]
All three statements are checked in Lean.
\end{theorem}
\begin{proof}
Represent one visible word by its finite sequence of spent bits and opaque saved tails followed by a live word. A positive spent slot surrounds exactly one such editable segment, with the unedited right duplicate opaque; a zero slot has no editable segment. A completed frame is the left context wrapping a packed list of these spent slots around its terminal \(S\) and a pristine allocation stamp. A well-formed outer list is one editable header followed by completed frames. In an open period, a second list of spent slots ends at the still-pristine suffix of the program. These are the inductive cases \texttt{SlotOK}, \texttt{FrameOK}, \texttt{OuterOK} and \texttt{Valid}.

For a word with spent prefix \(h\), the pop search either fails on an empty live suffix or gives the term with its first bit newly spent. Induct on \(|h|\), rebuilding the left continuation at each return. In a positive slot, rebuild the box, the open arm and the slot. In a zero slot take the next branch. Induction on the slot list and then on the outer frames shows that popping the bare term gives the same tree as erasing the annotated pop. The same inductions show the observable head agrees with the first live letter of the physical collection. Popping preserves its grammar: it changes only the indicated segment, and the copied history stays opaque. This also proves actual word readback including its chronological spent prefix.

A fire search skips all completed slots. If an open period is present, it therefore reaches its first pristine slot. The two arms give precisely the two terms of \eqref{commit}, and the remaining slots keep their order. At a frontier, firing fails, and allocation obeys \eqref{reset} or \eqref{alloc}. For the zero case the next frontier is \(X_{j+1}X_{j+1}\); the following allocation succeeds. For the positive case allocation succeeds immediately. The newly exposed period is well-formed. Firing either extends the open list by one spent slot or moves the completed period to the outer list. Thus every case preserves the grammar, and the unannotated action and annotated advancement return equal contraction counts and identical whole ground trees. Induction on \(n\) proves their iterates have equal erasures, including absorption at an empty word.

Independently, induction on each raw traversal shows that any success is a root \(S\) contraction in a one-hole context, even for inputs outside the grammar. The three possible successful branches have lengths \(2,3,4\). The compiler \(C\) first translates the transition table and initial tape into a primitive three-counter program and its initial registers. It next compiles that job to a well-formed restricted deletion-two tag job. It encodes the job for the Rogozhin interpreter and the fixed period-912 Cook system. The literal functions are \texttt{compileThreeCounterJob}, \texttt{compileT2Job} and \texttt{encodeBits} in Filin's \texttt{DeterministicTapeCook} module. The checked equivalences at these three interfaces compose to identify source halting with eventual cyclic dataword emptiness. These compilers terminate without using the source's halting time. Composing this theorem with the preceding identities proves the displayed equivalence. The formal theorem depends only on \texttt{propext} and \texttt{Quot.sound}.
\end{proof}

\begin{theorem}\label{microhalt}
Let \(\Delta\) be one call of the eight-mode controller, extended by the identity on normal terms. Put \(J(m,t)=[m=\mathsf{boundary}]H(t)\). For every deterministic tape-machine instance \(M\),
\[
 M\downarrow\quad\Longleftrightarrow\quad
 \exists k:\ J\!\left(\Delta^k(\mathsf{boundary},E(P_*,C(M)))\right)=1 .
\]
The guard may be tested after every individual contraction. It is a fixed regular test with one finite mode observation.
\end{theorem}
\begin{proof}
In a successful block of length \(l=2,3,4\), the interior modes are respectively the first \(l-1\) search stages with the same borrowed bit. Its endpoint is the boundary mode with exactly the tree returned by \(B\). Thus no interior stage accepts. Induction on the macro index realizes any first successful macro observation at some finite primitive time. Conversely, at any generated boundary with \(H=0\), the physical queue is nonempty and the next block exists. If a primitive-time acceptance occurs \(k\) steps later, either \(k=0\) and the initial boundary accepts, or \(k\) is at least the next block length \(l\). In the latter case shift to the resulting boundary with \(k-l<k\). Strong induction gives an accepting macro observation. An already accepting boundary needs no successor condition. Theorem~\ref{certmain} now gives the result. The complete argument is the declaration \texttt{UniformS.Micro.universal\_single\_contractions}.
\end{proof}

\paragraph{\textit{Finite output.}}
For \(d\) source machines, encode the \(d\) seeds in a fixed list of \(S\)-pairs. One round advances each branch once, rebuilding its unaltered siblings. Context closure gives at most \(4d\) native contractions per round. The \(i\)th detector is the preceding regular language at the \(i\)th list field. By induction its first success exists precisely when machine \(i\) halts. For two recognizers with disjoint halting domains, accept when either detector succeeds and return its index. The output is the literal identity of the accepting port. Every partial computable Boolean function admits this standard presentation: run its computing machine and accept precisely when its result is the designated bit, with a permanent loop on the opposite result. The two recognizers are finite programs, fixed independently of their input. The two-port statement is also formalized in Lean. This result differs from the single numerical menu of Theorem~\ref{output}, which returns an unbounded unary count from one physical queue.

\paragraph{\textit{Encoding and control costs.}}
Set \(m(0)=3,\ m(1)=4\), extending \(m\) additively to words. Unfolding \eqref{seed} and using structural induction gives the following kernel-checked exact leaf count:
\[
 L(E(P,w))=22+48p+4\sum_{i<p}m(v_i)+m(w).
\]
The bounded protocol can start each constituent search at the root. It passes only a bit and the finite stage between these searches; nothing is retained between macro calls. The bottom-up detector uses five bits and fifteen ternary values per summary. Direct evaluation with a recursion stack needs \(O(h)\) such summaries for a height-\(h\) occurrence tree. Each search is a linear-time scan with a bounded number of visits to each relevant occurrence. Parent links support root-restarted navigation as in Proposition~\ref{walker}. The direct Lean evaluator rebuilds persistent context, so its call stack and constructed term also consume space. These costs differ from the bounded number of contractions. On replacement systems of Theorem~\ref{cutoff} the same regular halt observations, even with feedback, yield a bounded abstract orbit. The concrete implementation admits no such replacement for its universal halting contract.

\section{Natural-number output on the binary evaluator}\label{numericv2}
Use the ordinary deterministic binary tape model: a finite table \(M\) assigns to a state and scanned bit at most one triple (write bit, head displacement in \(\{-1,0,1\}\), next state). A missing rule halts. A finite window grows by one blank on crossing either end. The initial state is zero, and the head starts at index one in \(0\,1^x0\). If \(r=(q,h,T)\) is terminal, put \(\nu(r)=|\operatorname{takeWhile}_{1}(\operatorname{drop}_{h}T)|\). This is the usual unary-output convention. Every finite binary tape table has a code \(e\): the paired lists of optional triples are numbered by the mutually inverse \texttt{machineCode} and \texttt{machineDecode} in the compiler of~\cite{filin}. Consequently this choice ranges over all deterministic Turing computations of partial numerical functions. A machine can blank the cell immediately after its output and position its head on its first unit. On the zero numeral it positions the head on a blank. These are ordinary source-machine conventions. They introduce no target-side oracle. Write \(U_T(e,x)\downarrow=y\) for reachability of a terminal row of value \(y\). Its numbering is independent of the register numbering of Theorem~\ref{output}.

Let \(P_*\) be the fixed 912-phase program in the table below, \(s=(M,0,1^x)\) and \(w=C(\operatorname{pad}s)\) the tape-to-Cook encoder of Cinematic Strawberry~\cite{filin}. Padding ensures that every finite tape table meets the source-side conditions. Its correctness in both directions is included in the checked dependency closure. Let \(B\) be the certified bare boundary policy, \(h\) its fixed regular empty-queue detector, and \(t_0(w)\) its previously defined unannotated seed. Set
\[
 E_T(e,x)=\pair{\word(w)}{t_0(w)},\qquad
 \sigma_T(\pair a t)=\pair a{B(t)},\qquad H_T(\pair a t)=h(t).
\]
On a nonmatching tree the policy fixes the tree and the guard is false. The extra copy of the seed is immutable. The constructor at its root and the right subtree are enough to recognize the guard by a finite tree automaton. All selections on a generated input start in the right subtree. For \(m(w)=3|w|_0+4|w|_1\), the exact unfolded encoder leaf count is
\(24+48|P_*|+4\sum_{v\in P_*}m(v)+2m(w)\).
This bound is in the compiled word length. Constructing \(w\) is effective, potentially enormous and separate from the listed contraction cost.

\paragraph{\textit{Reading physical data.}}
At a detected boundary let \(g\) be the bits of the visible cells of the right tree, in physical word order, dropping only their live/spent flags. All are spent at that boundary. Our reader uses the word length and bit pattern of the immutable copy; it receives no prior snapshot or externally supplied clock. Indices below start at zero. For suffix \(d\) and width \(u\), the following scanner returns a list of indexed queues:
\begin{align*}
 W(j,i,u,\varepsilon)&=[(j,i,\varepsilon)],\\
 W(j,i,u,bd)&=(j,i,\operatorname{take}_{u}(bd))::
 W\bigl(j+1,i+1\bmod |P_*|,(u\mathbin{\dot-}1)+b|v_i|,d\bigr).
\end{align*}
Here \(\dot-\) truncates at zero. The scanner reads at most \(|g|+1\) physical slices. Each bit delivered by a slice is present in \(g\). No transition produces a missing bit inside this decoder.

\begin{lemma}\label{fifoexact}
Suppose a run of a FIFO program \(P\) drains \(c=(i,z)\), consuming precisely \(g\). The corresponding scanner with inputs \(j,i,|z|,g\) returns exactly the sequence of actual configurations from \(c\) through the first empty queue, with consecutive indices starting at \(j\).
\end{lemma}
\begin{proof}
Represent complete consumptions inductively by \((i,\varepsilon;\varepsilon)\) and by prefixing \(b\) to the consumption of \((i+1,z'v_i^b)\) to obtain a consumption of \((i,bz')\). Induction in this relation first gives \(g=z d\) for some \(d\): a FIFO removes all of its initial suffix before its new appendant. Hence the first slice is exactly \(z\). When \(z=bz'\), the next width is \(|z'|+b|v_i|\), the phase increments, and the remaining suffix is precisely the consumption for that next queue. Apply induction. The empty case returns the single empty queue. For a generated term the native pop records the next actual bit and the native release appends exactly the appropriate word. Induction using the proved equality of the bare orbit to the erased board orbit yields precisely this complete-consumption relation whenever the regular detector fires. A stationary empty tail adds no physical cells.
\end{proof}

\paragraph{\textit{Literal tape parser.}}
For a slice \((j,i,z)\), form the partial row \(R_w(j,i,z)\) with the terminal-row reader \texttt{decodeTerminalTape?} of~\cite{filin}. Its data inputs are the static seed and the actual bits \(z\). At phase zero it reads aligned one-hot blocks of 114 bits. It uses the canonical decoder at index zero and the arrival decoder at a positive index. It parses the tagged boundary, deletion-two encoding and counter word, decodes the tape boundary, and tests the terminal rule in the seed-derived table. At other phases it fails. Each stage is a total option-valued reader. The positive index is used only to distinguish a later arrival from the initial word. Its magnitude is never used to run a machine for that many steps. The imported source-reflection theorem states, for each particular row \(r\),
\[
 \exists j:\ R_w(j,c_j)=r
 \quad\Longleftrightarrow\quad
 (\exists n:\operatorname{run}_{M}^{n}(r_0)=r)\ \land\ \operatorname{step}_M(r)=\bot ,
\]
where \(c_j\) is the actual \(j\)-th Cook configuration from \(w\). The kernel checks the full compiler, reflection and converse proof at the pinned revision. In addition, \(R_w(j,i,\varepsilon)=\bot\) for every \(i,j,w\): the empty one-hot sequence yields no canonical or arrival configuration. This last observation is essential because the indexed reflection theorem has no upper bound on \(j\).

Define \(D_T=\bot\) off the regular guard. At a hit apply \(R_w\) to the physical windows \(W(0,0,|w|,g)\), map each returned row through \(\nu\), and take the first successful value, returning \(\bot\) if none succeeds. Invalid finite trees have a total interpretation by the same option-valued structural definitions. Neither the encoder nor the decoder calls the source transition function, a run enumerator or an oracle for the result. The reflection argument invokes source execution only in its specification and proof.

\begin{theorem}\label{certnumeric}
For all \(e,x,y\in\mathbb N\),
\[
 U_T(e,x)\downarrow=y\quad\Longleftrightarrow\quad
 \exists k:\ H_T(\sigma_T^k E_T(e,x))=1,\quad D_T(\sigma_T^k E_T(e,x))=y .
\]
Every successful boundary transition uses two to four contextual \(S\) contractions. The same numerical equivalence holds when the eight-mode controller is sampled after every individual contraction. All accepted results for a given input agree.
\end{theorem}
\begin{proof}
By the kernel-checked regular domain equivalence, the guard is reached exactly on a halting input. At any hit, Lemma~\ref{fifoexact} makes membership in the list of physical readings equivalent to the existence of a successful reading of some \(c_j\): after its first empty queue every further \(c_j\) is empty and the empty reader always fails. The source-reflection equivalence now identifies the set of results with the unary values of the terminal rows of \(M\). A deterministic run reaching two terminal rows must reach the same row: induction over the two run lengths strips simultaneous steps, and an additional step from a terminal row is impossible. The first successful value therefore exists exactly for the correct \(y\). For the primitive controller, each bounded action has a block of exactly its declared length, with no intermediate boundary. On a productive generated seed every nonterminal boundary starts such a block. Conversely, induction on the number of primitive actions places each accepting primitive configuration at the identical bare tree of a boundary iterate. From an already accepting boundary the fallback either stutters or enters an absorbing nonaccepting mode. This proves both directions including all intermediate actions.
\end{proof}

\paragraph{\textit{Numerical separation.}}
The accepted boundary is stationary, so the decoder can be applied once, at the first guard hit. Consider any \(r\)-register replacement of Theorem~\ref{cutoff} with \(q\) possible summaries per register. Even granting an arbitrary total computable numerical reader of the \emph{entire concrete register vector} at its first accepting tick, its numerical input--output graph is decidable for every total computable encoder. Indeed if a first hit occurs at \(n\), the abstract cutoff supplies a hit \(j<q^r\). The concrete and abstract guards agree at every time, so firstness forces \(n\le j\). One can construct the concrete terms for fewer than \(q^r\) ticks and read the concrete value at the first hit. The Lean theorem \texttt{first\_numeric\_cutoff} proves this bounded equivalence for every reader. An enumeration with a verifier deciding \emph{which} time to accept does not have this restriction. Our output parser is invoked only after a regular native-dependent event whose eventual occurrence already characterizes the halting domain. The argument applies to the stated register replacements, including whole-tree copying and feedback. It asserts no limitation on arbitrary alternative active rewriting machines.

\section{Computational attribution}\label{attribution}
The two original observational constructions use the equivalence
\[
 F(e,x,c,z)\ \land\ \bigwedge_{j<b}Q_j=PT_j
 \quad\Longleftrightarrow\quad
 \exists i\,[z=z_i(e,x)\ \land\ c=c_i(e,x)] .
\]
The variable displacement between packed rows matters. The manuscript gives a concrete false halt if this comparison is dropped. Its Theorem 10.2 proves decidability on the specified classes of complete replacement records. Proposition 11.1 constructs the first checkpoint from an ordered unary clock by a fixed dimension-three first-order interpretation. For the second observer, the later repository note gives an exact quadratic-size shared graph and an observer-equivalent \emph{ordinary} tree of linear size produced in logspace. Its four readouts agree at every checkpoint \cite{observed,replacement}. The exact unfolded second checkpoint can be exponential. These comparisons preserve the history theorem and explain why its source separation leaves the attribution objection open. Checking a correct history is a mathematical computation; the issue here is that independent clock expansions supply all inputs needed by the same verifier.

In Theorem~\ref{main}, each nonempty source transition is realized from the preceding queue. The program is an arbitrary input to a single selector. The selector conveys one bit to the next slot. The last consumed bit determines which arm to fire. Its other decisions find the next cell, decide between pop and commit, and request allocation. The local walker skips the interior of the dormant appendant when selecting either arm. For every word \(v\), the same arm contraction exposes a whole residual of \(\word(v)\). Later pops rewrite that occurrence. The only state carried across invocations is the \(S\)-tree. The numerical specialization likewise obtains its next control label and output word from a physical record; the detector reads one finite label and parity, and the decoder counts one kind of living letter. It can equivalently emit one unary symbol per such occurrence.

\begin{prop}\label{clock}
For a fixed numerical template, a clock that leaves the header untouched and supplies an arbitrary input-independent rail cannot yield the universal halting domain with \emph{any} regular tree detector. The same conclusion holds for any regular detector after a fixed finite-copying MSO word-to-tree transduction of a standard unary input-and-tick word.
\end{prop}
\begin{proof}
On \(\pair{\bx{\word(R(\nu(e,x)))}}{C_n}\), any bottom-up detector sees the header only through one of finitely many states. Whether it can accept at some \(n\) is a fixed subset of these states, even if the sequence \(C_n\) is arbitrary. Compute the header state and test membership in that finite subset. This is a decidable input language. No uniform algorithm for extracting the subset from an arbitrary clock is claimed. For the second assertion, a run of a tree automaton can be pulled back to the copies of word positions by finitely many monadic state sets. Transition consistency and acceptance are MSO word conditions. Their word language is regular, and deleting an existential final tick block preserves regularity. Either conclusion contradicts universality of the halting domain in Theorem~\ref{output}.
\end{proof}

Two earlier constructions admit sharper bounds. The first has effectively Presburger visit sets for regular tree detectors, and regular tests on the exact second checkpoints have regular unary timing \cite{generator,replacement}. Their original arithmetic tests exceed those bounds.

A second check isolates the contribution of native copying. Replace the contraction by any single combinator rule with the same left side whose right side contains only application and uses each of \(x,y,z\) at most once, with no new constant. Each such rewrite strictly decreases leaf count. Every selected run then reaches a normal form in fewer than the initial number of leaves. Its observed halt problem for the same effective encoder and bounded readouts is decidable by direct execution. This excludes all such linear replacements even with the choice of occurrence still available.

These results address passive enumeration and the loss of copying. A different \emph{active} medium with unbounded fresh records, writable progress marks and conditional release of stored words can implement the same queue algorithm. In particular, the controller participates in the simulated transition by selecting its arm. The universality theorem uses a fixed regular choice of redex that depends on the current term, as in the selected-path formulation \cite{announcement,prize}.

\section{A replacement theorem with feedback}
\label{feedback}
A replacement can have writable, unbounded data and still admit a uniform decision procedure for every input. The following result makes that assertion precise. It also specifies the operation unavailable to this class: a polynomial update can copy an entire old register, but cannot expose one of its arbitrary interior subtrees as a new register.

Let \(Q\) be any finite algebra for a constant \(s\) and a binary operation \(\mu\). Write \(h:\T\to Q\) for its unique homomorphism. Taking products allows any fixed finite family of regular tree observations simultaneously. Choose \(r\) tree registers, an arbitrary function \(d:Q^r\to \T(x_1,\dots,x_r)^r\) and \(A\subseteq Q^r\). At a tick, evaluate the expressions \(d(h^r(\mathbf t))\) simultaneously in the old registers. Variables may occur any finite number of times. The register values can grow exponentially in the number of ticks even from a single small seed. Each fixed machine has a bounded per-tick growth factor. Acceptance reads \(A\) after any number of ticks, including zero. Initial registers may be computed by any total effective encoder.

\begin{theorem}\label{cutoff}
Write \(q=|Q|\). A replacement reaches \(A\) if and only if it does so at some \(k<q^r\). The induced halting domain is decidable. With a regular-preimage encoder, the halting domain on the original inputs is regular.
\end{theorem}
\begin{proof}
For each expression \(e\), structural induction gives
\[
 h(e_{\T}(\mathbf t))=e_Q(h^r(\mathbf t)).
\]
For a variable this is projection, for a constant it is the value \(s\), and for an application the induction hypotheses compose with \(\mu\). Thus the vector observation after one concrete tick is
\(\delta(a)=d(a)_Q(a)\), where \(a=h^r(\mathbf t)\). Induction gives equality after every number of ticks, even though the choice of expression depends on every preceding observation. If an earliest accepting index were at least \(q^r\), its prefix including that index would have two equal abstract states. Removing the segment between these states preserves the suffix up to acceptance, contradicting minimality. Enumerate the bounded orbit to decide. For the language claim, the acceptable initial vectors form a finite subset of \(Q^r\), so pull that subset back along the encoder.
\end{proof}

The bound concerns a joint observation state for an individual register. A finite external controller of size \(m\) can be included by taking a product on the vector orbit, giving \(m q^r\). A decoder may additionally count designated output occurrences after acceptance. As long as acceptance alone is equivalent to source halting, an unbounded numerical output cannot invalidate this decision procedure. The theorem does not include arbitrary rewriting at an observed address. In our construction a slot is opened and the designated continuation of its new word becomes the subject of later updates. On the one-word program \((1)\) with seed \(1\), every pop after the first is at the residual exposed by the immediately preceding firing. Induction on these firings proves this for all times. The position of the active residual descends through frames of unbounded depth. A clock presenting candidate numbers to an arithmetic history verifier has no such successor invariant. Applying the cutoff to such a clock requires a finite joint observation algebra for its verifier. The old arithmetic acceptance test has no such representation on the clock inputs that would give a decidable universal domain.

\begin{prop}\label{stack}
A finite controller on one unbounded pushdown stack, with regular whole-stack guards and a regular halting test, has decidable existential halting from every effectively constructed finite initial configuration. This also holds if guarded moves are nondeterministic.
\end{prop}
\begin{proof}
Use a product DFA for the guards and the target test, reading from the bottom. Decorate each stack symbol with the DFA state for the suffix below it. The state for the whole stack is obtained from the top and its decoration. A pop exposes the previous summary. A push installs the current summary below its new symbol. An endmarker handles emptiness. Finite top replacements maintain valid decorations by finitely many moves, giving an ordinary pushdown system with a regular target.

For completeness, form a finite automaton for that target with the pushdown controls as allowed starting states. Copy its internal automaton into disjoint states and take no initial edges entering controls. For each rule \(pA\to q u\) and existing path \(q\xrightarrow{u}v\), insert \(p\xrightarrow{A}v\). This uses finitely many possible edges, hence terminates. At the fixed point each actual predecessor of a recognized configuration is recognized: decompose its successor's accepting path after the prefix \(u\). Induction on execution length gives completeness. For soundness, in a saturated accepting path from a control, a first edge from the original automaton starts a wholly original accepting path since no added edge leaves an internal state. Otherwise replace its first, added edge by the earlier edges that justified it, and take the associated pushdown move. This strictly decreases the multiset of addition ranks along the accepting path, even if the replacement path is longer. Eventually the remaining accepting path is original, including the possible empty path. Thus the saturated automaton recognizes exactly the predecessor set. Test the initial decorated stack in it \cite{pushdown}.
\end{proof}

The two replacement theorems grant different operations. The first permits copying whole trees; the second allows popping and inspecting nested symbols. Neither covers a selector coupled to arbitrary occurrence-specific copying rewrites of a binary tree. These bounds exclude clock verification and both register and stack replacements with the stated observations. In the selected reduction, the walker carries the consumed bit to the next slot. Within this layout, moving that routing entirely into a fixed sequence confined to the slot is impossible when the stored word is nonempty. The two popped-bit cases have identical unused slots. Any identical sequence of addressed operations internal to those slots produces identical released payloads, whereas the two required outputs are \(\varepsilon\) and \(v\). A construction that joins the consumed cell and the slot in a new redex would fall outside this locality argument.

\section{Size, navigation and extraction}
Let \(N=2L-1\) denote unfolded nodes for a tree with \(L\) leaves; let \(h\) be its edge height. The direct encoder prints a constant number of program templates and one initial word. For input presented as a list of binary words, it takes \(O(1+p+|w|+\sum_i |v_i|)\) output nodes and time. Constructors add a bounded number of nodes per slot and letter. A finite-copying MSO interpretation suffices on a delimited input word: give each letter and delimiter a fixed finite set of copies for its local constructor, link consecutive cell and slot roots within their respective list, and attach the distinguishable endpoints to the constant seed nodes. Endpoints and next delimiters are first-order tests on the ordered input; fixed copies provide each occurrence of the program. This is an output-construction argument, not an invocation of the source interpreter. For a fixed numerical template the bound is \(O(1+e+x)\) with unary inputs. It is exponential in the worst case if \(e,x\) are supplied in binary. Printing the constant numerical menu has a substantial one-time cost. No such linear construction bound is claimed for converting an arbitrary Turing-machine description to the upstream fixed Cook input.

With parent and incoming-side navigation on the unfolded occurrence tree, the walker has finite control, one read-only cursor and no writable stack. A machine storing the cursor as an index into an explicit table needs \(O(\log(N+1))\) bits for that index; parent links and incoming sides belong to the input representation. With downward-only immutable references the executable implements a zipper using \(O(h)\) references for the current context. Storing a path as a repeatedly concatenated string can cost quadratic time in its length. A strict finite tree automaton for the detector uses constant information per node and a postorder implementation needs \(O(h)\) stack frames. The numerical decoder emits one symbol per designated live unit using the parsed roles; a binary count needs \(O(\log(N+1))\) result bits. Even if implemented by an explicit stack or temporary annotations, a scan reads only the current finite term. The operations of the model are independent of such implementation choices.

For a contraction with third argument \(z\), \(L'=L-1+L(z)\). The local height increases by at most one. These facts persist in a surrounding context (take sums for leaves and maxima for height). Induction gives \(L_k\le 2^k L_0\) and \(h_k\le h_0+k\). Hence repeated linear tree scans have total cost \(O((2^k-1)N_0+k)\) through \(k\) contractions, apart from initialization. Fresh copies of the full term need \(O(2^k N_0)\) peak tree space if old trees can be released. Immutable DAG evaluation may reuse the two references to \(z\). For an address of depth \(d\), allocate at most three application nodes for the right side and \(d\) copied ancestors. Unfolding this rebuild of the unique occurrence gives the exact contextual tree replacement, by induction on the address. Even retaining every previous root, the number of allocated DAG nodes is at most
\[
 D_0+3k+k h_0+k(k-1)/2 .
\]
This is a storage bound for nodes constructed by contraction and path copying. It excludes selector workspaces and memo tables, and it gives no corresponding polynomial runtime for scanning the unfolded tree. The two-to-four count excludes navigation, printing the seed, conversion from another source model, the number of letters per register sweep, and cost of materializing an output.

\section{Computational checks and related work}\label{sources}
The reference implementation contracts redexes in immutable application trees, rebuilding ancestors so only the selected occurrence changes. It compares every stable queue with an independent source trajectory and checks the intermediate parity difference after every pop or allocation. In the source-level register comparisons, a separate implementation of increment, decrement and jump checks the complete numeric row. Cached facts in the optional full-interpreter run are deterministic functions of immutable menu subtrees. They accelerate reconstruction of the selected occurrence. The local walker of Proposition~\ref{walker} reconstructs every nomination from scratch.

\begin{center}
\small
\begin{tabular}{@{}lr@{}}
\toprule
Test & Result\\
\midrule
Cyclic program/input seeds & 916\\
Source steps taking 2, 3, and 4 contractions & 3,286;\ 1,662;\ 2,150\\
Finite-queue transitions in direct register comparisons & 1,008\\
Register instructions checked on 30 seeds & 171\\
Native contractions across the bounded regression suite & 25,763\\
Fresh-root choices compared to the regular scan & 27,480\\
All bare trees with at most nine leaves checked & 2,056\\
Complete interpreter run \(U(1,0)=0\) & 6,552 contractions\\
Complete interpreter run \(U(1,1)=1\) & 11,668 contractions\\
\bottomrule
\end{tabular}
\end{center}
The first 840 cyclic cases exhaust all periods one and two, appendant lengths at most two, and seed lengths at most three, up to twelve source steps or empty termination. Another 75 programs are deterministically sampled and one singleton program follows 128 successive bits, including 127 released by earlier steps. The two full runs use object code \(1\), the program that immediately returns its input, on the actual 446-label interpreter. It executes 66 and 108 primitive register instructions, respectively. The separate root-walk replay compares 30 consecutive nominations per cyclic seed and checks legality or normality on every bare tree through nine leaves. The physical-window test checks a further 224 completed FIFO runs. For each run, the windows read from the final tree agree with the queues recorded by an independent tag interpreter.

\paragraph{\textit{Word readers and program allocation.}}
An earlier observational construction encoded computation histories in selected terms \cite{observed}. Its clock replacements are discussed in Section~\ref{attribution} \cite{replacement}. A subsequent register simulation used a regular selector and an allocator for writable slots \cite{five}. Here each newly allocated frame holds an entire period of the cyclic program. One contraction releases a word of arbitrary finite length, and subsequent contractions consume that copy in queue order.

Langdon's word reader distinguishes live and spent cells and admits a finite controller that restarts at the root \cite{langdonword}. Equation~\eqref{read} uses this convention. A related selected-port identity illustrates the role of the contraction address \cite{langdonport}. For \(H_C(D)=S(SC)D\), contraction at the root followed by two contractions at its right child gives
\[
 S C H_C(D)Y \longrightarrow_S^3 (CY)((C(DY))(Y(DY))).
\]
In our simulation, equations~\eqref{reset} and~\eqref{alloc} supply the next program period when the current one has been exhausted.

\paragraph{\textit{Earlier universality under selected reduction.}}
Cinematic Strawberry, in work directed by Alexander Filin, established pure-\(S\) universality under a fixed root-restarted controller with a regular halting language \cite{filin}. Its Theorem~1R gives cyclic tag trajectories and ordered source observations. That construction specializes a dispatcher to the cyclic program. In the present construction the same controller acts on every finite program, which is stored in the initial term. Our numerical proof uses Cinematic Strawberry's tape-to-tag compiler, numbering of machine tables, and terminal-row reader. The FIFO representation and the proof that the reader's input occurs in the native term are given here.

\section{Formal verification}
In Lean, \(\mathrm{Native}(t,u)\) denotes one contextual \(S\) contraction, and \(\mathrm{Steps}(k,t,u)\) a chain of exactly \(k\) contractions. The auxiliary type \(\mathrm{Board}\) describes the visible segments, spent cells, committed slots and frontier of a term. Erasing a board \(b\) gives an ordinary closed tree \(\lfloor b\rfloor_P\). The boundary policy of Theorem~\ref{certmain} acts directly on that tree. The stateless selector of Theorem~\ref{main} is represented separately. On any finite term its fallback finds a redex if one exists.

Write \(b_n=\mathrm{sample}(P,n,b_0)\) for the deterministic composition of the annotated pop, allocation and firing operations, with a stationary sample after an empty queue. Theorems \texttt{macro\_exact} and \texttt{ground\_trace} prove the 2--4 bounds and the source state equation for every nonempty transition, arbitrary program, and arbitrary finite nonempty-step run, including all previously copied history in the ground-tree conclusion. The spent prefix has length equal to the number of nonempty source steps. A second induction establishes the following theorem, checked in \texttt{PureS/Tape.lean}. Fix Filin's Rogozhin--Cook cyclic program \(P_*\). For every deterministic tape-machine instance \(M\), let \(b_0\) be the seed of its compiled input. Then
\[
 M\downarrow\quad\Longleftrightarrow\quad
 \exists n,k:\ k\le4n,\quad
 \lfloor b_0\rfloor_{P_*}\xrightarrow{k}_{S}\lfloor b_n\rfloor_{P_*},
 \quad V(b_n)=\varepsilon .
\]
The sequence \(b_n\) in this equivalence is determined by the encoded machine and its input. The tape-to-Cook compiler and its halting theorem are taken from \cite{filin}. The complete certificate includes 90 modules from that development, together with the proofs relating its cyclic queues and terminal rows to the present \(S\)-terms.

\paragraph{\textit{Scope of the formalization.}}
In \texttt{UniformS.Numerical}, Theorem~\ref{certnumeric} takes two forms. At tag-step boundaries it is \nolinkurl{universal_natural_io}. After each single contraction it is \nolinkurl{universal_single_natural_io}. Both depend on propositional extensionality and quotient soundness. The first-hit bound for tree registers also uses classical choice. The two numerical declarations and their dependencies passed Lean~4.33.1's fresh-kernel check. A second checker, \texttt{lean4lean}, replayed all 112 source modules.

The proof of Theorem~\ref{main} identifies the address chosen by the stateless parity walker on each reachable term. Its Lean development proves legality, exact reporting of normality, and agreement on spent words and initial seeds. The full correspondence for that walker is proved in Section~\ref{regular}. For the register convention of Theorem~\ref{output}, Lean checks a complete primitive sweep on every legal row. The finite-menu simulation of the 446-label interpreter and its register-4 readout are proved in the paper. Theorem~\ref{certnumeric} uses a different machine numbering and has an end-to-end Lean proof.

\paragraph{\textit{A common test instance.}}
For \(P=(1,\varepsilon)\) and initial word \(101\), both our construction and that of \cite{filin} visit the queues \(101,011,11\). Here the corresponding contraction indices are \(0,4,6\), and the unfolded trees have \(289,635,643\) nodes. The earlier construction reports indices \(0,22,85\) and tree sizes \(171,17057,339285\). These measurements refer to the same two source steps. The general bound in Theorem~\ref{main} concerns the number of contractions.

The \href{https://github.com/Bobtheotherone/pure-s-universality}{supplementary material} includes the Lean sources, the reference evaluator, a Wolfram Language example, and instructions for reproducing the proofs and computations.


\section*{Acknowledgments}
OpenAI Codex assisted with the mathematical derivations, Lean formalization, reference implementation, testing, and preparation of the manuscript.

\appendix
\section{The universal numerical interpreter}\label{app:machine}
Here is a self-contained choice of the indexing and interpreter used in Theorem~\ref{output}. It fixes the target independently of the selected reduction. An object program is a finite list of instructions on numbered natural-valued registers. An increment of \(j\) jumps to label \(l\). A conditional decrement of \(j\) jumps to \(l_0\) on zero and to \(l_1\) after a positive decrement. Halt returns register zero. Initially only register zero can be nonzero, with value \(x\). An unoccupied instruction address or a malformed opcode enters an infinite loop. An unconditional jump in an object program uses a permanently zero spare register.

Represent nil by zero and a list cell with head \(a\) and tail \(b\) by
\[
       [a\mid b]=2^a(2b+1).
\]
Every positive integer splits uniquely, and its tail is smaller, so lists are finite. Instruction zero is halt. The increment instruction has code \([0\mid[j\mid l]]\). The conditional decrement has code \([1\mid[j\mid[l_0\mid l_1]]]\). The outer program list of instruction codes has natural code \(e\), with label zero at its head. Require the inner cells in a nonhalting instruction to be positive. The unit program consisting of halt has code \(1\). The program that increments register zero, jumps to label one, and halts has code \([7\mid 1]=384\).

\paragraph{\textit{Interpreter.}}
Reserve two untouched registers \(E,X\) and initialize separate working registers to their values. All other primitive registers start at zero. Set the virtual register list to \([x\mid0]\) and the virtual program counter to zero. The following fixed finite control can be read independently of the accompanying executable.
\begin{enumerate}
\item Copy the program list to a scan register and drop as many cells as specified by the virtual program counter, rejecting a missing instruction by entering a self-jump. Split the selected cell to obtain its instruction head.
\item For halt, take the head of the virtual register list into the designated primitive output register and enter the absorbing halt label. For another instruction, split its fields, check its tag and required inner cells, and retain \(j\) and the branch target or targets.
\item Traverse the virtual register list to index \(j\), saving skipped heads on a reversed temporary list. For register access only, splitting nil gives head and tail zero. This implements implicit trailing zero registers. Split off the addressed value, preserving its tail.
\item Execute the specified operation on this value and set the next virtual program counter. Cons the new value onto the retained tail. Restore saved heads by popping the reversed temporary list and consing them in turn. Return to instruction fetch.
\end{enumerate}

These operations admit a finite primitive expansion. Clear a register by guarded decrements. Transfer its units with a decrement/increment loop. To copy without losing the source, increment the destination and a zero temporary for each consumed unit, then restore the source from that temporary. Double a value with two increments per unit, using an empty temporary. To cons, start with twice the tail plus one and double once per unit of a working copy of the head. To split a positive number, remove pairs of units to form its half and record the parity in finite control. At an even division increment the prospective head and repeat on the quotient. At the first odd division the current quotient is the tail. Even positive quotients strictly decrease before another division. Scratch values at macro boundaries are cleared or consumed. The protected registers \(E,X\) are never decremented, even transiently.

Induction on the virtual program steps shows that a fetch sees exactly the virtual machine's program counter and registers. In a finite access, the traversals are bounded by the supplied address or by the finite reversed prefix. All arithmetic subroutines on a finite input terminate. Thus one defined virtual step takes finitely many primitive steps, the error loop diverges, and a virtual halt returns the designated output. The particular expansion in the companion has 29 registers (including temporaries) and 446 labeled primitive instructions. Its initialization uses only the zero and copy wiring of equation~\eqref{maps}. A choice of different hygienic temporaries or an additional jump would change the target's microstep trace. The queue simulation applies at each of its primitive instructions.

\paragraph{\textit{Universality of the indexing.}}
For completeness, a fixed Turing machine can be compiled into a finite object program. Keep the state and scanned tape symbol in finite control and store the strings strictly left and strictly right of the head in two integer stacks, nearest symbol first. With an alphabet encoded by nonzero digits \(1,\ldots,B-1\) for fixed \(B\), let \(b_\ast\) denote the digit for a blank. Push digit \(d\) as \(Bs+d\); pop by quotient and remainder modulo \(B\). On an empty stack return blank and leave it empty. Blank is also an allowed stored digit. For a right move, push the written symbol on the left stack and pop the right stack into the scanned symbol. A left move uses the symmetric operations. A fixed-\(B\) product is repeated addition. Quotient and remainder are obtained by subtracting units while cycling through \(B\) finite control labels and incrementing a quotient on each completed cycle. Temporary registers implement copies and restores as above. This proves one-step agreement with the Turing machine by induction, with only finitely many object instructions per machine description. A unary input tape is initialized by decrementing a copy of \(x\), and a Turing machine computing a numeric function may return its result as a delimited unary block, counted into register zero. This gives all partial computable functions. Conversely the interpreter's primitive instructions are effective. These facts give the stated universal partial function and its undecidable halting domain, in agreement with classical register-machine universality \cite{register}.

\Needspace{13\baselineskip}
\section{Primitive code of the numerical interpreter}\label{app:code}
The following table gives every entry in address order. The first member of a tuple encodes the operation. The remaining members are operands. For program counter \(q\), the instructions are
\[
 \begin{gathered}(0):\mathrm{halt},\qquad (1,i,h):(n_i,q)\gets(n_i+1,h),\qquad(3,h):q\gets h,\\
 (2,i,h,j): (n_i,q)\gets
 \begin{cases}(0,h)&n_i=0,\\(n_i-1,j)&n_i>0.\end{cases}\end{gathered}
\]
Input slots are \((e,x,e,x)\) followed by zeros. A halt returns slot 4.
{\fontsize{9}{11}\selectfont
\setlength{\tabcolsep}{4pt}
\begin{longtable}{@{}rl@{\qquad}rl@{\qquad}rl@{}}
\toprule
\(q\)&\(I_q\)&\(q\)&\(I_q\)&\(q\)&\(I_q\)\\
\midrule
\endhead
\bottomrule
\endlastfoot
\(0\)&\((2,5,1,0)\) & \(1\)&\((2,3,5,2)\) & \(2\)&\((1,5,3)\)\\
\(3\)&\((1,6,4)\) & \(4\)&\((3,1)\) & \(5\)&\((2,6,8,6)\)\\
\(6\)&\((1,3,7)\) & \(7\)&\((3,5)\) & \(8\)&\((2,7,9,8)\)\\
\(9\)&\((2,8,13,10)\) & \(10\)&\((1,7,11)\) & \(11\)&\((1,6,12)\)\\
\(12\)&\((3,9)\) & \(13\)&\((2,6,16,14)\) & \(14\)&\((1,8,15)\)\\
\(15\)&\((3,13)\) & \(16\)&\((2,7,20,17)\) & \(17\)&\((1,9,18)\)\\
\(18\)&\((1,9,19)\) & \(19\)&\((3,16)\) & \(20\)&\((2,9,23,21)\)\\
\(21\)&\((1,7,22)\) & \(22\)&\((3,20)\) & \(23\)&\((1,7,24)\)\\
\(24\)&\((2,5,33,25)\) & \(25\)&\((2,7,29,26)\) & \(26\)&\((1,9,27)\)\\
\(27\)&\((1,9,28)\) & \(28\)&\((3,25)\) & \(29\)&\((2,9,32,30)\)\\
\(30\)&\((1,7,31)\) & \(31\)&\((3,29)\) & \(32\)&\((3,24)\)\\
\(33\)&\((2,10,34,33)\) & \(34\)&\((2,7,37,35)\) & \(35\)&\((1,10,36)\)\\
\(36\)&\((3,34)\) & \(37\)&\((2,11,38,37)\) & \(38\)&\((2,2,42,39)\)\\
\(39\)&\((1,11,40)\) & \(40\)&\((1,6,41)\) & \(41\)&\((3,38)\)\\
\(42\)&\((2,6,45,43)\) & \(43\)&\((1,2,44)\) & \(44\)&\((3,42)\)\\
\(45\)&\((2,12,46,45)\) & \(46\)&\((2,13,50,47)\) & \(47\)&\((1,12,48)\)\\
\(48\)&\((1,6,49)\) & \(49\)&\((3,46)\) & \(50\)&\((2,6,53,51)\)\\
\(51\)&\((1,13,52)\) & \(52\)&\((3,50)\) & \(53\)&\((2,11,420,54)\)\\
\(54\)&\((1,11,55)\) & \(55\)&\((2,12,81,56)\) & \(56\)&\((2,14,57,56)\)\\
\(57\)&\((2,11,61,58)\) & \(58\)&\((1,14,59)\) & \(59\)&\((1,6,60)\)\\
\(60\)&\((3,57)\) & \(61\)&\((2,6,64,62)\) & \(62\)&\((1,11,63)\)\\
\(63\)&\((3,61)\) & \(64\)&\((2,15,65,64)\) & \(65\)&\((2,11,66,65)\)\\
\(66\)&\((2,14,80,67)\) & \(67\)&\((1,14,68)\) & \(68\)&\((2,14,72,69)\)\\
\(69\)&\((2,14,77,70)\) & \(70\)&\((1,16,71)\) & \(71\)&\((3,68)\)\\
\(72\)&\((2,16,75,73)\) & \(73\)&\((1,14,74)\) & \(74\)&\((3,72)\)\\
\(75\)&\((1,15,76)\) & \(76\)&\((3,68)\) & \(77\)&\((2,16,80,78)\)\\
\(78\)&\((1,11,79)\) & \(79\)&\((3,77)\) & \(80\)&\((3,53)\)\\
\(81\)&\((2,14,82,81)\) & \(82\)&\((2,11,86,83)\) & \(83\)&\((1,14,84)\)\\
\(84\)&\((1,6,85)\) & \(85\)&\((3,82)\) & \(86\)&\((2,6,89,87)\)\\
\(87\)&\((1,11,88)\) & \(88\)&\((3,86)\) & \(89\)&\((2,17,90,89)\)\\
\(90\)&\((2,15,91,90)\) & \(91\)&\((2,14,105,92)\) & \(92\)&\((1,14,93)\)\\
\(93\)&\((2,14,97,94)\) & \(94\)&\((2,14,102,95)\) & \(95\)&\((1,16,96)\)\\
\(96\)&\((3,93)\) & \(97\)&\((2,16,100,98)\) & \(98\)&\((1,14,99)\)\\
\(99\)&\((3,97)\) & \(100\)&\((1,17,101)\) & \(101\)&\((3,93)\)\\
\(102\)&\((2,16,105,103)\) & \(103\)&\((1,15,104)\) & \(104\)&\((3,102)\)\\
\(105\)&\((2,17,421,106)\) & \(106\)&\((1,17,107)\) & \(107\)&\((2,14,108,107)\)\\
\(108\)&\((2,17,112,109)\) & \(109\)&\((1,14,110)\) & \(110\)&\((1,6,111)\)\\
\(111\)&\((3,108)\) & \(112\)&\((2,6,115,113)\) & \(113\)&\((1,17,114)\)\\
\(114\)&\((3,112)\) & \(115\)&\((2,18,116,115)\) & \(116\)&\((2,19,117,116)\)\\
\(117\)&\((2,14,131,118)\) & \(118\)&\((1,14,119)\) & \(119\)&\((2,14,123,120)\)\\
\(120\)&\((2,14,128,121)\) & \(121\)&\((1,16,122)\) & \(122\)&\((3,119)\)\\
\(123\)&\((2,16,126,124)\) & \(124\)&\((1,14,125)\) & \(125\)&\((3,123)\)\\
\(126\)&\((1,18,127)\) & \(127\)&\((3,119)\) & \(128\)&\((2,16,131,129)\)\\
\(129\)&\((1,19,130)\) & \(130\)&\((3,128)\) & \(131\)&\((2,19,420,132)\)\\
\(132\)&\((1,19,133)\) & \(133\)&\((2,14,134,133)\) & \(134\)&\((2,19,138,135)\)\\
\(135\)&\((1,14,136)\) & \(136\)&\((1,6,137)\) & \(137\)&\((3,134)\)\\
\(138\)&\((2,6,141,139)\) & \(139\)&\((1,19,140)\) & \(140\)&\((3,138)\)\\
\(141\)&\((2,20,142,141)\) & \(142\)&\((2,21,143,142)\) & \(143\)&\((2,14,157,144)\)\\
\(144\)&\((1,14,145)\) & \(145\)&\((2,14,149,146)\) & \(146\)&\((2,14,154,147)\)\\
\(147\)&\((1,16,148)\) & \(148\)&\((3,145)\) & \(149\)&\((2,16,152,150)\)\\
\(150\)&\((1,14,151)\) & \(151\)&\((3,149)\) & \(152\)&\((1,20,153)\)\\
\(153\)&\((3,145)\) & \(154\)&\((2,16,157,155)\) & \(155\)&\((1,21,156)\)\\
\(156\)&\((3,154)\) & \(157\)&\((2,22,158,157)\) & \(158\)&\((2,18,189,159)\)\\
\(159\)&\((2,21,420,160)\) & \(160\)&\((1,21,161)\) & \(161\)&\((2,18,164,162)\)\\
\(162\)&\((1,18,163)\) & \(163\)&\((3,420)\) & \(164\)&\((2,14,165,164)\)\\
\(165\)&\((2,21,169,166)\) & \(166\)&\((1,14,167)\) & \(167\)&\((1,6,168)\)\\
\(168\)&\((3,165)\) & \(169\)&\((2,6,172,170)\) & \(170\)&\((1,21,171)\)\\
\(171\)&\((3,169)\) & \(172\)&\((2,23,173,172)\) & \(173\)&\((2,24,174,173)\)\\
\(174\)&\((2,14,188,175)\) & \(175\)&\((1,14,176)\) & \(176\)&\((2,14,180,177)\)\\
\(177\)&\((2,14,185,178)\) & \(178\)&\((1,16,179)\) & \(179\)&\((3,176)\)\\
\(180\)&\((2,16,183,181)\) & \(181\)&\((1,14,182)\) & \(182\)&\((3,180)\)\\
\(183\)&\((1,23,184)\) & \(184\)&\((3,176)\) & \(185\)&\((2,16,188,186)\)\\
\(186\)&\((1,24,187)\) & \(187\)&\((3,185)\) & \(188\)&\((1,22,189)\)\\
\(189\)&\((2,11,190,189)\) & \(190\)&\((2,10,194,191)\) & \(191\)&\((1,11,192)\)\\
\(192\)&\((1,6,193)\) & \(193\)&\((3,190)\) & \(194\)&\((2,6,197,195)\)\\
\(195\)&\((1,10,196)\) & \(196\)&\((3,194)\) & \(197\)&\((2,25,198,197)\)\\
\(198\)&\((2,20,261,199)\) & \(199\)&\((2,14,200,199)\) & \(200\)&\((2,11,204,201)\)\\
\(201\)&\((1,14,202)\) & \(202\)&\((1,6,203)\) & \(203\)&\((3,200)\)\\
\(204\)&\((2,6,207,205)\) & \(205\)&\((1,11,206)\) & \(206\)&\((3,204)\)\\
\(207\)&\((2,26,208,207)\) & \(208\)&\((2,11,209,208)\) & \(209\)&\((2,14,223,210)\)\\
\(210\)&\((1,14,211)\) & \(211\)&\((2,14,215,212)\) & \(212\)&\((2,14,220,213)\)\\
\(213\)&\((1,16,214)\) & \(214\)&\((3,211)\) & \(215\)&\((2,16,218,216)\)\\
\(216\)&\((1,14,217)\) & \(217\)&\((3,215)\) & \(218\)&\((1,26,219)\)\\
\(219\)&\((3,211)\) & \(220\)&\((2,16,223,221)\) & \(221\)&\((1,11,222)\)\\
\(222\)&\((3,220)\) & \(223\)&\((2,5,224,223)\) & \(224\)&\((2,26,228,225)\)\\
\(225\)&\((1,5,226)\) & \(226\)&\((1,6,227)\) & \(227\)&\((3,224)\)\\
\(228\)&\((2,6,231,229)\) & \(229\)&\((1,26,230)\) & \(230\)&\((3,228)\)\\
\(231\)&\((2,7,232,231)\) & \(232\)&\((2,25,236,233)\) & \(233\)&\((1,7,234)\)\\
\(234\)&\((1,6,235)\) & \(235\)&\((3,232)\) & \(236\)&\((2,6,239,237)\)\\
\(237\)&\((1,25,238)\) & \(238\)&\((3,236)\) & \(239\)&\((2,7,243,240)\)\\
\(240\)&\((1,9,241)\) & \(241\)&\((1,9,242)\) & \(242\)&\((3,239)\)\\
\(243\)&\((2,9,246,244)\) & \(244\)&\((1,7,245)\) & \(245\)&\((3,243)\)\\
\(246\)&\((1,7,247)\) & \(247\)&\((2,5,256,248)\) & \(248\)&\((2,7,252,249)\)\\
\(249\)&\((1,9,250)\) & \(250\)&\((1,9,251)\) & \(251\)&\((3,248)\)\\
\(252\)&\((2,9,255,253)\) & \(253\)&\((1,7,254)\) & \(254\)&\((3,252)\)\\
\(255\)&\((3,247)\) & \(256\)&\((2,25,257,256)\) & \(257\)&\((2,7,260,258)\)\\
\(258\)&\((1,25,259)\) & \(259\)&\((3,257)\) & \(260\)&\((3,198)\)\\
\(261\)&\((2,14,262,261)\) & \(262\)&\((2,11,266,263)\) & \(263\)&\((1,14,264)\)\\
\(264\)&\((1,6,265)\) & \(265\)&\((3,262)\) & \(266\)&\((2,6,269,267)\)\\
\(267\)&\((1,11,268)\) & \(268\)&\((3,266)\) & \(269\)&\((2,27,270,269)\)\\
\(270\)&\((2,28,271,270)\) & \(271\)&\((2,14,285,272)\) & \(272\)&\((1,14,273)\)\\
\(273\)&\((2,14,277,274)\) & \(274\)&\((2,14,282,275)\) & \(275\)&\((1,16,276)\)\\
\(276\)&\((3,273)\) & \(277\)&\((2,16,280,278)\) & \(278\)&\((1,14,279)\)\\
\(279\)&\((3,277)\) & \(280\)&\((1,27,281)\) & \(281\)&\((3,273)\)\\
\(282\)&\((2,16,285,283)\) & \(283\)&\((1,28,284)\) & \(284\)&\((3,282)\)\\
\(285\)&\((2,22,305,286)\) & \(286\)&\((2,27,296,287)\) & \(287\)&\((2,13,288,287)\)\\
\(288\)&\((2,24,292,289)\) & \(289\)&\((1,13,290)\) & \(290\)&\((1,6,291)\)\\
\(291\)&\((3,288)\) & \(292\)&\((2,6,295,293)\) & \(293\)&\((1,24,294)\)\\
\(294\)&\((3,292)\) & \(295\)&\((3,314)\) & \(296\)&\((2,13,297,296)\)\\
\(297\)&\((2,23,301,298)\) & \(298\)&\((1,13,299)\) & \(299\)&\((1,6,300)\)\\
\(300\)&\((3,297)\) & \(301\)&\((2,6,304,302)\) & \(302\)&\((1,23,303)\)\\
\(303\)&\((3,301)\) & \(304\)&\((3,314)\) & \(305\)&\((1,27,306)\)\\
\(306\)&\((2,13,307,306)\) & \(307\)&\((2,21,311,308)\) & \(308\)&\((1,13,309)\)\\
\(309\)&\((1,6,310)\) & \(310\)&\((3,307)\) & \(311\)&\((2,6,314,312)\)\\
\(312\)&\((1,21,313)\) & \(313\)&\((3,311)\) & \(314\)&\((2,5,315,314)\)\\
\(315\)&\((2,27,319,316)\) & \(316\)&\((1,5,317)\) & \(317\)&\((1,6,318)\)\\
\(318\)&\((3,315)\) & \(319\)&\((2,6,322,320)\) & \(320\)&\((1,27,321)\)\\
\(321\)&\((3,319)\) & \(322\)&\((2,7,323,322)\) & \(323\)&\((2,28,327,324)\)\\
\(324\)&\((1,7,325)\) & \(325\)&\((1,6,326)\) & \(326\)&\((3,323)\)\\
\(327\)&\((2,6,330,328)\) & \(328\)&\((1,28,329)\) & \(329\)&\((3,327)\)\\
\(330\)&\((2,7,334,331)\) & \(331\)&\((1,9,332)\) & \(332\)&\((1,9,333)\)\\
\(333\)&\((3,330)\) & \(334\)&\((2,9,337,335)\) & \(335\)&\((1,7,336)\)\\
\(336\)&\((3,334)\) & \(337\)&\((1,7,338)\) & \(338\)&\((2,5,347,339)\)\\
\(339\)&\((2,7,343,340)\) & \(340\)&\((1,9,341)\) & \(341\)&\((1,9,342)\)\\
\(342\)&\((3,339)\) & \(343\)&\((2,9,346,344)\) & \(344\)&\((1,7,345)\)\\
\(345\)&\((3,343)\) & \(346\)&\((3,338)\) & \(347\)&\((2,11,348,347)\)\\
\(348\)&\((2,7,351,349)\) & \(349\)&\((1,11,350)\) & \(350\)&\((3,348)\)\\
\(351\)&\((2,25,415,352)\) & \(352\)&\((1,25,353)\) & \(353\)&\((2,14,354,353)\)\\
\(354\)&\((2,25,358,355)\) & \(355\)&\((1,14,356)\) & \(356\)&\((1,6,357)\)\\
\(357\)&\((3,354)\) & \(358\)&\((2,6,361,359)\) & \(359\)&\((1,25,360)\)\\
\(360\)&\((3,358)\) & \(361\)&\((2,26,362,361)\) & \(362\)&\((2,25,363,362)\)\\
\(363\)&\((2,14,377,364)\) & \(364\)&\((1,14,365)\) & \(365\)&\((2,14,369,366)\)\\
\(366\)&\((2,14,374,367)\) & \(367\)&\((1,16,368)\) & \(368\)&\((3,365)\)\\
\(369\)&\((2,16,372,370)\) & \(370\)&\((1,14,371)\) & \(371\)&\((3,369)\)\\
\(372\)&\((1,26,373)\) & \(373\)&\((3,365)\) & \(374\)&\((2,16,377,375)\)\\
\(375\)&\((1,25,376)\) & \(376\)&\((3,374)\) & \(377\)&\((2,5,378,377)\)\\
\(378\)&\((2,26,382,379)\) & \(379\)&\((1,5,380)\) & \(380\)&\((1,6,381)\)\\
\(381\)&\((3,378)\) & \(382\)&\((2,6,385,383)\) & \(383\)&\((1,26,384)\)\\
\(384\)&\((3,382)\) & \(385\)&\((2,7,386,385)\) & \(386\)&\((2,11,390,387)\)\\
\(387\)&\((1,7,388)\) & \(388\)&\((1,6,389)\) & \(389\)&\((3,386)\)\\
\(390\)&\((2,6,393,391)\) & \(391\)&\((1,11,392)\) & \(392\)&\((3,390)\)\\
\(393\)&\((2,7,397,394)\) & \(394\)&\((1,9,395)\) & \(395\)&\((1,9,396)\)\\
\(396\)&\((3,393)\) & \(397\)&\((2,9,400,398)\) & \(398\)&\((1,7,399)\)\\
\(399\)&\((3,397)\) & \(400\)&\((1,7,401)\) & \(401\)&\((2,5,410,402)\)\\
\(402\)&\((2,7,406,403)\) & \(403\)&\((1,9,404)\) & \(404\)&\((1,9,405)\)\\
\(405\)&\((3,402)\) & \(406\)&\((2,9,409,407)\) & \(407\)&\((1,7,408)\)\\
\(408\)&\((3,406)\) & \(409\)&\((3,401)\) & \(410\)&\((2,11,411,410)\)\\
\(411\)&\((2,7,414,412)\) & \(412\)&\((1,11,413)\) & \(413\)&\((3,411)\)\\
\(414\)&\((3,351)\) & \(415\)&\((2,10,416,415)\) & \(416\)&\((2,11,419,417)\)\\
\(417\)&\((1,10,418)\) & \(418\)&\((3,416)\) & \(419\)&\((3,37)\)\\
\(420\)&\((3,420)\) & \(421\)&\((2,14,422,421)\) & \(422\)&\((2,10,426,423)\)\\
\(423\)&\((1,14,424)\) & \(424\)&\((1,6,425)\) & \(425\)&\((3,422)\)\\
\(426\)&\((2,6,429,427)\) & \(427\)&\((1,10,428)\) & \(428\)&\((3,426)\)\\
\(429\)&\((2,4,430,429)\) & \(430\)&\((2,15,431,430)\) & \(431\)&\((2,14,445,432)\)\\
\(432\)&\((1,14,433)\) & \(433\)&\((2,14,437,434)\) & \(434\)&\((2,14,442,435)\)\\
\(435\)&\((1,16,436)\) & \(436\)&\((3,433)\) & \(437\)&\((2,16,440,438)\)\\
\(438\)&\((1,14,439)\) & \(439\)&\((3,437)\) & \(440\)&\((1,4,441)\)\\
\(441\)&\((3,433)\) & \(442\)&\((2,16,445,443)\) & \(443\)&\((1,15,444)\)\\
\(444\)&\((3,442)\) & \(445\)&\((0)\) & &\\
\end{longtable}}

\Needspace{13\baselineskip}
\section{The fixed cyclic program}
The fixed program has 912 phases. The table lists the first 114 productions as words of symbol indices. A superscript denotes repetition. The remaining 798 appendants are empty. The following substitution gives the binary word appended at each phase. The tape-to-word translation is that of~\cite{filin}.
\[
 A_j=0^j1\,0^{113-j},\quad
 v_i=\begin{cases}\prod_{j\in \tau_i}^{\prec}A_j&i<114,\\
 \varepsilon&114\le i<912,\end{cases}\quad
 \langle j\rangle^k=\underbrace{(j)\cdot(j)\cdots(j)}_{k}.
\]
{\fontsize{9}{11}\selectfont
\setlength{\tabcolsep}{5pt}
\begin{longtable}{@{}llll@{}}
\toprule
\endfirsthead
\toprule
\endhead
\bottomrule
\endlastfoot
\multicolumn{4}{@{}l@{}}{\(0\mapsto \langle 16\rangle^{1}\langle 17\rangle^{1}\langle 18\rangle^{1}\langle 19\rangle^{1}\langle 20\rangle^{1}\langle 21\rangle^{1}\langle 22\rangle^{1}\langle 23\rangle^{1}\)}\\
\multicolumn{4}{@{}l@{}}{\(1\mapsto \langle 24\rangle^{1}\langle 25\rangle^{1}\langle 26\rangle^{1}\langle 27\rangle^{1}\langle 28\rangle^{1}\langle 29\rangle^{1}\langle 30\rangle^{1}\langle 31\rangle^{1}\)}\\
\multicolumn{4}{@{}l@{}}{\(2\mapsto \langle 32\rangle^{1}\langle 33\rangle^{1}\langle 34\rangle^{1}\langle 35\rangle^{1}\langle 36\rangle^{1}\langle 37\rangle^{1}\langle 38\rangle^{1}\langle 39\rangle^{1}\)}\\
\multicolumn{4}{@{}l@{}}{\(3\mapsto \langle 40\rangle^{1}\langle 41\rangle^{1}\langle 42\rangle^{1}\langle 43\rangle^{1}\langle 44\rangle^{1}\langle 45\rangle^{1}\langle 46\rangle^{1}\langle 47\rangle^{1}\)}\\
\multicolumn{4}{@{}l@{}}{\(4\mapsto \langle 48\rangle^{1}\langle 49\rangle^{1}\langle 50\rangle^{1}\langle 51\rangle^{1}\langle 52\rangle^{1}\langle 53\rangle^{1}\langle 54\rangle^{1}\langle 55\rangle^{1}\)}\\
\multicolumn{4}{@{}l@{}}{\(5\mapsto \langle 56\rangle^{1}\langle 57\rangle^{1}\langle 58\rangle^{1}\langle 59\rangle^{1}\langle 60\rangle^{1}\langle 61\rangle^{1}\langle 62\rangle^{1}\langle 63\rangle^{1}\)}\\
\multicolumn{4}{@{}l@{}}{\(6\mapsto \langle 64\rangle^{1}\langle 65\rangle^{1}\langle 66\rangle^{1}\langle 67\rangle^{1}\langle 68\rangle^{1}\langle 69\rangle^{1}\langle 70\rangle^{1}\langle 71\rangle^{1}\)}\\
\multicolumn{4}{@{}l@{}}{\(7\mapsto \langle 72\rangle^{1}\langle 73\rangle^{1}\langle 74\rangle^{1}\langle 75\rangle^{1}\langle 76\rangle^{1}\langle 77\rangle^{1}\langle 78\rangle^{1}\langle 79\rangle^{1}\)}\\
\multicolumn{4}{@{}l@{}}{\(8\mapsto \langle 80\rangle^{1}\langle 81\rangle^{1}\langle 82\rangle^{1}\langle 83\rangle^{1}\langle 84\rangle^{1}\langle 85\rangle^{1}\langle 86\rangle^{1}\langle 87\rangle^{1}\)}\\
\multicolumn{4}{@{}l@{}}{\(9\mapsto \langle 88\rangle^{1}\langle 89\rangle^{1}\langle 90\rangle^{1}\langle 91\rangle^{1}\langle 92\rangle^{1}\langle 93\rangle^{1}\langle 94\rangle^{1}\langle 95\rangle^{1}\)}\\
\multicolumn{4}{@{}l@{}}{\(10\mapsto \langle 96\rangle^{1}\langle 97\rangle^{1}\langle 98\rangle^{1}\langle 99\rangle^{1}\langle 100\rangle^{1}\langle 101\rangle^{1}\langle 102\rangle^{1}\langle 103\rangle^{1}\)}\\
\multicolumn{4}{@{}l@{}}{\(11\mapsto \langle 104\rangle^{1}\langle 105\rangle^{1}\langle 106\rangle^{1}\langle 107\rangle^{1}\langle 108\rangle^{1}\langle 109\rangle^{1}\langle 110\rangle^{1}\langle 111\rangle^{1}\)}\\
\midrule
\(12\mapsto \langle 8\rangle^{8}\) & \(13\mapsto \langle 9\rangle^{8}\) & \(14\mapsto \langle 10\rangle^{8}\) & \(15\mapsto \langle 11\rangle^{8}\)\\
\(16\mapsto \langle 12\rangle^{32}\langle 0\rangle^{1}\) & \(17\mapsto \langle 0\rangle^{2}\langle 4\rangle^{40}\) & \(18\mapsto \langle 12\rangle^{48}\langle 0\rangle^{3}\) & \(19\mapsto \langle 0\rangle^{4}\langle 4\rangle^{24}\)\\
\(20\mapsto \langle 12\rangle^{32}\langle 0\rangle^{5}\) & \(21\mapsto \langle 3\rangle^{6}\langle 7\rangle^{24}\) & \(22\mapsto \langle 0\rangle^{10}\langle 4\rangle^{8}\) & \(23\mapsto \langle 0\rangle^{11}\)\\
\(24\mapsto \langle 1\rangle^{1}\langle 5\rangle^{24}\) & \(25\mapsto \langle 14\rangle^{40}\langle 2\rangle^{2}\) & \(26\mapsto \langle 1\rangle^{3}\langle 5\rangle^{32}\) & \(27\mapsto \langle 13\rangle^{40}\langle 1\rangle^{4}\)\\
\(28\mapsto \langle 13\rangle^{56}\langle 1\rangle^{5}\) & \(29\mapsto \langle 1\rangle^{6}\langle 5\rangle^{48}\) & \(30\mapsto \langle 1\rangle^{10}\langle 5\rangle^{8}\) & \(31\mapsto \langle 1\rangle^{11}\)\\
\(32\mapsto \langle 2\rangle^{1}\langle 6\rangle^{56}\) & \(33\mapsto \langle 3\rangle^{2}\langle 7\rangle^{32}\) & \(34\mapsto \langle 2\rangle^{3}\langle 6\rangle^{48}\) & \(35\mapsto \varepsilon\)\\
\(36\mapsto \langle 0\rangle^{5}\langle 4\rangle^{16}\) & \(37\mapsto \langle 0\rangle^{6}\langle 4\rangle^{56}\) & \(38\mapsto \langle 2\rangle^{10}\langle 6\rangle^{8}\) & \(39\mapsto \langle 2\rangle^{11}\)\\
\(40\mapsto \langle 3\rangle^{1}\langle 7\rangle^{24}\) & \(41\mapsto \langle 13\rangle^{16}\langle 1\rangle^{2}\) & \(42\mapsto \langle 3\rangle^{3}\langle 7\rangle^{32}\) & \(43\mapsto \varepsilon\)\\
\(44\mapsto \langle 13\rangle^{16}\langle 1\rangle^{5}\) & \(45\mapsto \langle 3\rangle^{6}\langle 7\rangle^{48}\) & \(46\mapsto \langle 3\rangle^{10}\langle 7\rangle^{8}\) & \(47\mapsto \langle 3\rangle^{11}\)\\
\(48\mapsto \langle 4\rangle^{1}\) & \(49\mapsto \langle 4\rangle^{64}\) & \(50\mapsto \langle 4\rangle^{1}\) & \(51\mapsto \langle 4\rangle^{64}\)\\
\(52\mapsto \langle 4\rangle^{1}\) & \(53\mapsto \langle 7\rangle^{64}\) & \(54\mapsto \langle 4\rangle^{8}\) & \(55\mapsto \langle 4\rangle^{8}\)\\
\(56\mapsto \langle 5\rangle^{64}\) & \(57\mapsto \langle 6\rangle^{1}\) & \(58\mapsto \langle 5\rangle^{64}\) & \(59\mapsto \langle 5\rangle^{1}\)\\
\(60\mapsto \langle 5\rangle^{1}\) & \(61\mapsto \langle 5\rangle^{64}\) & \(62\mapsto \langle 5\rangle^{8}\) & \(63\mapsto \langle 5\rangle^{8}\)\\
\(64\mapsto \langle 6\rangle^{64}\) & \(65\mapsto \langle 7\rangle^{64}\) & \(66\mapsto \langle 6\rangle^{64}\) & \(67\mapsto \varepsilon\)\\
\(68\mapsto \langle 4\rangle^{64}\) & \(69\mapsto \langle 4\rangle^{64}\) & \(70\mapsto \langle 6\rangle^{8}\) & \(71\mapsto \langle 6\rangle^{8}\)\\
\(72\mapsto \langle 7\rangle^{64}\) & \(73\mapsto \langle 5\rangle^{1}\) & \(74\mapsto \langle 7\rangle^{64}\) & \(75\mapsto \varepsilon\)\\
\(76\mapsto \langle 5\rangle^{1}\) & \(77\mapsto \langle 7\rangle^{64}\) & \(78\mapsto \langle 7\rangle^{8}\) & \(79\mapsto \langle 7\rangle^{8}\)\\
\(80\mapsto \langle 8\rangle^{64}\) & \(81\mapsto \langle 8\rangle^{1}\) & \(82\mapsto \langle 8\rangle^{64}\) & \(83\mapsto \langle 8\rangle^{1}\)\\
\(84\mapsto \langle 8\rangle^{64}\) & \(85\mapsto \langle 11\rangle^{1}\) & \(86\mapsto \langle 8\rangle^{8}\) & \(87\mapsto \langle 8\rangle^{8}\)\\
\(88\mapsto \langle 9\rangle^{1}\) & \(89\mapsto \langle 10\rangle^{64}\) & \(90\mapsto \langle 9\rangle^{1}\) & \(91\mapsto \langle 9\rangle^{64}\)\\
\(92\mapsto \langle 9\rangle^{64}\) & \(93\mapsto \langle 9\rangle^{1}\) & \(94\mapsto \langle 9\rangle^{8}\) & \(95\mapsto \langle 9\rangle^{8}\)\\
\(96\mapsto \langle 10\rangle^{1}\) & \(97\mapsto \langle 11\rangle^{1}\) & \(98\mapsto \langle 10\rangle^{1}\) & \(99\mapsto \varepsilon\)\\
\(100\mapsto \langle 8\rangle^{1}\) & \(101\mapsto \langle 8\rangle^{1}\) & \(102\mapsto \langle 10\rangle^{8}\) & \(103\mapsto \langle 10\rangle^{8}\)\\
\(104\mapsto \langle 11\rangle^{1}\) & \(105\mapsto \langle 9\rangle^{64}\) & \(106\mapsto \langle 11\rangle^{1}\) & \(107\mapsto \varepsilon\)\\
\(108\mapsto \langle 9\rangle^{64}\) & \(109\mapsto \langle 11\rangle^{1}\) & \(110\mapsto \langle 11\rangle^{8}\) & \(111\mapsto \langle 11\rangle^{8}\)\\
\(112\mapsto \varepsilon\) & \(113\mapsto \varepsilon\) &  & \\
\end{longtable}}

\clearpage
\begingroup\raggedright
\bibliographystyle{plain}
\bibliography{references}
\endgroup
\end{document}
